% !PRINTED 235
\weberpage{235}

%<<CONT
\gesperrt{determinant} $-1$ \gesperrt{as substitutions of the second kind}.

According to the results of \S.\,50, these groups are isomorphic with the groups of binary linear substitutions with determinant $+1$ and $-1$, whereas for the fractional substitutions, which led us to the polyhedral groups, the two kinds are not distinguishable.

The finite groups of the second kind are, apart from their general group-theoretic interest, important on account of their application in crystallography. Their complete determination no longer presents any essential difficulty. By substitutions we understand simply binary linear substitutions with determinant $\pm 1$, and we recall that two such substitutions whose coefficients differ only in their common sign, such as
\[
\begin{pmatrix} \alpha, & \beta \\ \gamma, & \delta \end{pmatrix}, \quad
\begin{pmatrix} -\alpha, & -\beta \\ -\gamma, & -\delta \end{pmatrix},
\]
are regarded as identical (\S.\,49).

Since a group of linear substitutions contains the identical substitution, which is of determinant $+1$, in every group $Q$ of the second kind there must occur, besides the improper substitutions, proper ones as well, and these proper substitutions form by themselves a group $P$. If $\varphi$ is any substitution of the second kind occurring in $Q$, then every composition of $\varphi$ with a substitution of the second kind from $Q$ occurs in $P$, and hence we can always decompose $Q$ thus:
\weq{1}{Q = P + P\varphi.}
The group $P$ must be one of the polyhedral groups considered earlier. We must have $\varphi^{-1}P\varphi = P$, and $P$ is therefore a \gesperrt{normal divisor} of $Q$. Conversely, if $P$ is a polyhedral group and $\varphi$ a substitution of the second kind that satisfies the condition $\varphi^{-1}P\varphi = P$, then $Q = P + P\varphi$ is a \gesperrt{group of the second kind}.

Here too we regard two groups that arise from one another by transformation as not essentially different. In order to form all $Q$, we have to go through the various cases.

I. If $P$ is the unit group, which we denote as the cyclic group of the first degree by $C_1$, so that $P$ consists only of the
%CONT>>
